%% kapa: exam=1 round=28 year=2017 period=1 form=A subject=economics subject_ko=경제학원론 seq=37 number=77
%% source: Q-Net 감정평가사 기출문제 (cst00309) artlSeq=5199603 file=extracted/1008/5199603/2017년 제28회 감정평가사 1차시험 문제지/1교시 A형.pdf
%% conversion: kordoc 4.24.51 -> tools/gen1.py (auto); review=proofread-sonnet; flags=none
\begin{question}{77}
단기 총공급곡선에 관한 설명으로 옳은 것은?
\begin{choices}
  \item 케인즈(J. M. Keynes)에 따르면 명목임금이 고정되어 있는 단기에서 물가가 상승하면 고용량이 증가하여 생산량이 증가한다.
  \item 가격경직성 모형(sticky-price model)에서 물가수준이 기대 물가수준보다 낮다면 생산량은 자연산출량 수준보다 높다.
  \item 가격경직성 모형은 기업들이 가격수용자라고 전제한다.
  \item 불완전정보 모형(imperfect information model)은 가격에 대한 불완전한 정보로 인하여 시장은 불균형을 이룬다고 가정한다.
  \item 불완전정보 모형에서 기대 물가수준이 상승하면 단기 총공급곡선은 오른쪽으로 이동된다.
\end{choices}
\end{question}
\answer{1}
